\begin{deluxetable*}{ccccccccccccc}
\tablecolumns{13}
\tablewidth{0pt}
\tabletypesize{\scriptsize}
\tablecaption{\label{tab:stack}
Characteristics of the XXL samples}
\tablehead{
\multicolumn{1}{c}{Sample} & 
\multicolumn{1}{c}{$N_\mathrm{cl}$\tablenotemark{a}} & 
\multicolumn{1}{c}{$N_{T}$\tablenotemark{b}}  & 
\multicolumn{1}{c}{$T_\mathrm{X}$\tablenotemark{c}} & 
\multicolumn{1}{c}{$\langle T_\mathrm{X}\rangle_\mathrm{wl}$} & 
\multicolumn{1}{c}{$\zmed$\tablenotemark{d}} & 
\multicolumn{1}{c}{$\langle z\rangle_\mathrm{wl}$} &  
\multicolumn{1}{c}{$c_{200}$} & 
\multicolumn{1}{c}{$M_{200}$} & 
\multicolumn{1}{c}{$\langle M_{200}\rangle_\mathrm{wl}$} & 
\multicolumn{1}{c}{$\langle M_{200}\rangle_\mathrm{g}$} & 
\multicolumn{1}{c}{SNR} & 
\multicolumn{1}{c}{$($SNR$)_\mathrm{q}$} \\
\colhead{} & 
\colhead{} & 
\colhead{} & 
\multicolumn{1}{c}{(keV)} & 
\multicolumn{1}{c}{(keV)} & 
\colhead{} & 
\colhead{} & 
\colhead{} & 
\multicolumn{1}{c}{($10^{13}h^{-1}M_\odot$)} & 
\multicolumn{1}{c}{($10^{13}h^{-1}M_\odot$)} & 
\multicolumn{1}{c}{($10^{13}h^{-1}M_\odot$)} & 
\colhead{} & 
\colhead{}  
}
\startdata
C1+C2 & $136$ & $105$ & $1.9$ & $2.0$ & $0.30$ & $0.25$ & $3.5\pm 0.9$ & $8.7\pm 0.8$ & $8.0\pm 0.8$ & $9.8\pm 0.8$ & $15.6$ & $20.5$\\
C1 & $83$ & $76$ & $2.1$ & $2.1$ & $0.29$ & $0.23$ & $3.6\pm 1.1$ & $9.7\pm 1.0$ & $9.0\pm 1.0$ & $11.6\pm 1.2$ & $14.0$ & $18.4$\\
C2 & $53$ & $29$ & $1.7$ & $1.6$ & $0.43$ & $0.29$ & $3.4\pm 1.8$ & $6.4\pm 1.2$ & $6.1\pm 1.1$ & $6.5\pm 1.0$ & $7.2$ & $9.5$
\enddata
\tablenotetext{a}{Number of clusters.}
\tablenotetext{b}{Number of clusters with measured X-ray temperatures.}
\tablenotetext{c}{Median X-ray temperature.}
\tablenotetext{d}{Median cluster redshift.}
\tablecomments{Quantities in brackets with subscript "wl" denote lensing-weighted sample means (Equation (\ref{eq:wlmean})), and those in brackets with subscript "g"e denote error-weighted geometric means (Equation (\ref{eq:geom})). The effective mass and concentration parameters ($M_{200}, c_{200}$) of each subsample are obtained from a single-mass-bin fit to the respective stacked $\Delta\Sigma$ profile assuming an NFW density profile. For each sample, the effective $M_{200}$ mass is consistent with the respective weighted sample averages from individual cluster weak-lensing measurements. We provide two different estimates of the weak-lensing signal-to-nose ratio integrated over the comoving radial range $R\in [0.3, 3]\,\Mpch$, one based on the linear estimator, SNR (Equation (\ref{eq:SNR})), and the other based on the quadratic estimator, $($SNR$)_\mathrm{q}$ (Equation (\ref{eq:SNRq})).}
\end{deluxetable*}
